%%%%%%%%%%%%%%%%%%%%%%%%%%%%%%%%%%%%%%%%%%%%%%%%%%%%%%%%%%%%
% 8.1 DESCRIPTIVE STATISTICS
% The Composition of Advanced Mathematics
%%%%%%%%%%%%%%%%%%%%%%%%%%%%%%%%%%%%%%%%%%%%%%%%%%%%%%%%%%%%

\section{Descriptive Statistics}
\label{sec:descriptive-statistics}

\prerequisite{
Basic arithmetic, algebra, functions, summation notation, sets,
probability distributions, random variables, expectation, variance,
and introductory data interpretation.
}

\learningobjectives{
By the end of this section, the reader should be able to:
\begin{itemize}
    \item distinguish populations, samples, parameters, and statistics;
    \item distinguish categorical and quantitative variables;
    \item classify nominal, ordinal, discrete, and continuous data;
    \item organize data using frequency and relative-frequency tables;
    \item interpret bar charts, histograms, and empirical distributions;
    \item compute and interpret the arithmetic mean;
    \item compute weighted means;
    \item compute and interpret medians and modes;
    \item compare resistant and nonresistant measures of center;
    \item compute ranges, interquartile ranges, and deviations;
    \item compute population and sample variance;
    \item compute population and sample standard deviation;
    \item understand degrees of freedom in sample variance;
    \item standardize observations using \(z\)-scores;
    \item compute quartiles, percentiles, and quantiles;
    \item construct and interpret five-number summaries;
    \item identify possible outliers using the \(1.5\,IQR\) rule;
    \item interpret boxplots;
    \item compute moments and measures of skewness;
    \item interpret kurtosis cautiously;
    \item understand empirical distribution functions;
    \item compute covariance and correlation from data;
    \item distinguish association from causation;
    \item summarize grouped data;
    \item recognize the effect of transformations on descriptive
          statistics;
    \item understand robust statistics;
    \item compare multiple groups using standardized summaries;
    \item connect descriptive statistics with probability distributions,
          estimation, visualization, and statistical inference.
\end{itemize}
}

%%%%%%%%%%%%%%%%%%%%%%%%%%%%%%%%%%%%%%%%%%%%%%%%%%%%%%%%%%%%
\subsection{What Is Descriptive Statistics?}
\label{subsec:what-is-descriptive-statistics}
%%%%%%%%%%%%%%%%%%%%%%%%%%%%%%%%%%%%%%%%%%%%%%%%%%%%%%%%%%%%

Statistics is the mathematical study of data.

A typical statistical investigation begins with observed values such as

\[
x_1,x_2,\ldots,x_n.
\]

Descriptive statistics summarizes the data without yet making formal
probabilistic conclusions about a larger population.

Its goals include describing:

\[
\boxed{
\text{center},
}
\]

\[
\boxed{
\text{spread},
}
\]

\[
\boxed{
\text{shape},
}
\]

and

\[
\boxed{
\text{relationships}.
}
\]

Thus descriptive statistics transforms a raw collection of observations
into interpretable numerical and graphical summaries.

%%%%%%%%%%%%%%%%%%%%%%%%%%%%%%%%%%%%%%%%%%%%%%%%%%%%%%%%%%%%
\subsection{Population and Sample}
\label{subsec:population-sample}
%%%%%%%%%%%%%%%%%%%%%%%%%%%%%%%%%%%%%%%%%%%%%%%%%%%%%%%%%%%%

\begin{definitionbox}{Population}
A population is the complete collection of individuals, objects,
measurements, or outcomes of interest in a statistical investigation.
\end{definitionbox}

\begin{definitionbox}{Sample}
A sample is a subset of observations selected from the population.
\end{definitionbox}

If the population values are

\[
X_1,\ldots,X_N,
\]

then the population size is

\[
N.
\]

A sample may contain

\[
n
\]

observations, where usually

\[
n<N.
\]

%%%%%%%%%%%%%%%%%%%%%%%%%%%%%%%%%%%%%%%%%%%%%%%%%%%%%%%%%%%%
\subsection{Parameter and Statistic}
\label{subsec:parameter-statistic}
%%%%%%%%%%%%%%%%%%%%%%%%%%%%%%%%%%%%%%%%%%%%%%%%%%%%%%%%%%%%

\begin{definitionbox}{Parameter}
A parameter is a numerical characteristic of a population.
\end{definitionbox}

Examples include:

\[
\boxed{
\mu
=
\text{population mean},
}
\]

\[
\boxed{
\sigma^2
=
\text{population variance},
}
\]

and

\[
\boxed{
p
=
\text{population proportion}.
}
\]

\begin{definitionbox}{Statistic}
A statistic is a numerical quantity calculated from sample data.
\end{definitionbox}

Examples include:

\[
\boxed{
\overline{x}
=
\text{sample mean},
}
\]

\[
\boxed{
s^2
=
\text{sample variance},
}
\]

and

\[
\boxed{
\widehat{p}
=
\text{sample proportion}.
}
\]

%%%%%%%%%%%%%%%%%%%%%%%%%%%%%%%%%%%%%%%%%%%%%%%%%%%%%%%%%%%%
\subsection{Parameter Versus Statistic}
\label{subsec:parameter-versus-statistic}
%%%%%%%%%%%%%%%%%%%%%%%%%%%%%%%%%%%%%%%%%%%%%%%%%%%%%%%%%%%%

A useful distinction is

\[
\boxed{
\text{parameter}
=
\text{population quantity},
}
\]

while

\[
\boxed{
\text{statistic}
=
\text{sample quantity}.
}
\]

A parameter is typically fixed but unknown.

A statistic varies from sample to sample and is therefore a random
variable before the data are observed.

This distinction becomes fundamental in statistical inference.

%%%%%%%%%%%%%%%%%%%%%%%%%%%%%%%%%%%%%%%%%%%%%%%%%%%%%%%%%%%%
\subsection{Observations and Variables}
\label{subsec:observations-variables}
%%%%%%%%%%%%%%%%%%%%%%%%%%%%%%%%%%%%%%%%%%%%%%%%%%%%%%%%%%%%

An observation is one measured unit of data.

A variable is a characteristic recorded for each observation.

For example, a student dataset might contain variables such as:

\[
\text{major},
\]

\[
\text{age},
\]

\[
\text{exam score},
\]

and

\[
\text{credits completed}.
\]

Each row may represent one student.

Each column may represent one variable.

%%%%%%%%%%%%%%%%%%%%%%%%%%%%%%%%%%%%%%%%%%%%%%%%%%%%%%%%%%%%
\subsection{Categorical Variables}
\label{subsec:categorical-variables}
%%%%%%%%%%%%%%%%%%%%%%%%%%%%%%%%%%%%%%%%%%%%%%%%%%%%%%%%%%%%

\begin{definitionbox}{Categorical Variable}
A categorical variable records membership in a category rather than a
meaningful numerical magnitude.
\end{definitionbox}

Examples include:

\[
\text{blood type},
\]

\[
\text{department},
\]

\[
\text{color},
\]

and

\[
\text{transportation mode}.
\]

Categorical variables are sometimes called qualitative variables.

%%%%%%%%%%%%%%%%%%%%%%%%%%%%%%%%%%%%%%%%%%%%%%%%%%%%%%%%%%%%
\subsection{Nominal Variables}
\label{subsec:nominal-variables}
%%%%%%%%%%%%%%%%%%%%%%%%%%%%%%%%%%%%%%%%%%%%%%%%%%%%%%%%%%%%

A nominal variable consists of categories without a natural ordering.

Examples include:

\[
\text{red, blue, green},
\]

or

\[
\text{biology, mathematics, physics}.
\]

The numerical codes

\[
1,2,3
\]

may be assigned to categories for storage, but those codes do not imply
an arithmetic ordering.

%%%%%%%%%%%%%%%%%%%%%%%%%%%%%%%%%%%%%%%%%%%%%%%%%%%%%%%%%%%%
\subsection{Ordinal Variables}
\label{subsec:ordinal-variables}
%%%%%%%%%%%%%%%%%%%%%%%%%%%%%%%%%%%%%%%%%%%%%%%%%%%%%%%%%%%%

An ordinal variable has categories with a meaningful order.

Examples include:

\[
\text{poor},
\quad
\text{fair},
\quad
\text{good},
\quad
\text{excellent}.
\]

The categories can be ranked, but the numerical distance between ranks
need not be meaningful.

%%%%%%%%%%%%%%%%%%%%%%%%%%%%%%%%%%%%%%%%%%%%%%%%%%%%%%%%%%%%
\subsection{Quantitative Variables}
\label{subsec:quantitative-variables}
%%%%%%%%%%%%%%%%%%%%%%%%%%%%%%%%%%%%%%%%%%%%%%%%%%%%%%%%%%%%

\begin{definitionbox}{Quantitative Variable}
A quantitative variable records numerical measurements for which
arithmetic operations have meaningful interpretations.
\end{definitionbox}

Examples include:

\[
\text{height},
\]

\[
\text{mass},
\]

\[
\text{temperature},
\]

\[
\text{income},
\]

and

\[
\text{number of defects}.
\]

%%%%%%%%%%%%%%%%%%%%%%%%%%%%%%%%%%%%%%%%%%%%%%%%%%%%%%%%%%%%
\subsection{Discrete Variables}
\label{subsec:discrete-variables-statistics}
%%%%%%%%%%%%%%%%%%%%%%%%%%%%%%%%%%%%%%%%%%%%%%%%%%%%%%%%%%%%

A discrete quantitative variable takes values from a finite or countable
set.

Examples include:

\[
\text{number of customers},
\]

\[
\text{number of children},
\]

and

\[
\text{number of machine failures}.
\]

Typical values are integers.

%%%%%%%%%%%%%%%%%%%%%%%%%%%%%%%%%%%%%%%%%%%%%%%%%%%%%%%%%%%%
\subsection{Continuous Variables}
\label{subsec:continuous-variables-statistics}
%%%%%%%%%%%%%%%%%%%%%%%%%%%%%%%%%%%%%%%%%%%%%%%%%%%%%%%%%%%%

A continuous quantitative variable can theoretically take any value in
an interval.

Examples include:

\[
\text{time},
\]

\[
\text{distance},
\]

\[
\text{temperature},
\]

and

\[
\text{mass}.
\]

Observed values may be rounded because measurement instruments have
finite precision.

%%%%%%%%%%%%%%%%%%%%%%%%%%%%%%%%%%%%%%%%%%%%%%%%%%%%%%%%%%%%
\subsection{Levels of Measurement}
\label{subsec:levels-of-measurement}
%%%%%%%%%%%%%%%%%%%%%%%%%%%%%%%%%%%%%%%%%%%%%%%%%%%%%%%%%%%%

A traditional classification uses four measurement levels:

\[
\boxed{
\text{nominal},
}
\]

\[
\boxed{
\text{ordinal},
}
\]

\[
\boxed{
\text{interval},
}
\]

and

\[
\boxed{
\text{ratio}.
}
\]

Nominal and ordinal scales are categorical.

Interval and ratio scales are numerical.

%%%%%%%%%%%%%%%%%%%%%%%%%%%%%%%%%%%%%%%%%%%%%%%%%%%%%%%%%%%%
\subsection{Interval Scale}
\label{subsec:interval-scale}
%%%%%%%%%%%%%%%%%%%%%%%%%%%%%%%%%%%%%%%%%%%%%%%%%%%%%%%%%%%%

An interval scale has meaningful differences but no meaningful absolute
zero.

A common example is temperature measured in degrees Celsius.

The statement

\[
20^\circ\text{C}
-
10^\circ\text{C}
=
10^\circ\text{C}
\]

is meaningful.

But saying

\[
20^\circ\text{C}
\]

is twice as hot as

\[
10^\circ\text{C}
\]

is not meaningful in the ratio-scale sense.

%%%%%%%%%%%%%%%%%%%%%%%%%%%%%%%%%%%%%%%%%%%%%%%%%%%%%%%%%%%%
\subsection{Ratio Scale}
\label{subsec:ratio-scale}
%%%%%%%%%%%%%%%%%%%%%%%%%%%%%%%%%%%%%%%%%%%%%%%%%%%%%%%%%%%%

A ratio scale has meaningful differences and a meaningful zero.

Examples include:

\[
\text{length},
\]

\[
\text{mass},
\]

\[
\text{time duration},
\]

and

\[
\text{counts}.
\]

For ratio data, multiplicative comparisons can be meaningful.

%%%%%%%%%%%%%%%%%%%%%%%%%%%%%%%%%%%%%%%%%%%%%%%%%%%%%%%%%%%%
\subsection{Univariate, Bivariate, and Multivariate Data}
\label{subsec:univariate-bivariate-multivariate}
%%%%%%%%%%%%%%%%%%%%%%%%%%%%%%%%%%%%%%%%%%%%%%%%%%%%%%%%%%%%

Data involving one variable are univariate.

Data involving two variables are bivariate.

Data involving three or more variables are multivariate.

Thus:

\[
\boxed{
\text{univariate}
\rightarrow
\text{one variable},
}
\]

\[
\boxed{
\text{bivariate}
\rightarrow
\text{two variables},
}
\]

\[
\boxed{
\text{multivariate}
\rightarrow
\text{several variables}.
}
\]

%%%%%%%%%%%%%%%%%%%%%%%%%%%%%%%%%%%%%%%%%%%%%%%%%%%%%%%%%%%%
\subsection{Raw Data}
\label{subsec:raw-data}
%%%%%%%%%%%%%%%%%%%%%%%%%%%%%%%%%%%%%%%%%%%%%%%%%%%%%%%%%%%%

Suppose a sample consists of

\[
4,\ 7,\ 6,\ 4,\ 9,\ 7,\ 5,\ 4.
\]

These unsummarized observations are called raw data.

Descriptive statistics seeks to organize them into useful forms.

%%%%%%%%%%%%%%%%%%%%%%%%%%%%%%%%%%%%%%%%%%%%%%%%%%%%%%%%%%%%
\subsection{Frequency}
\label{subsec:frequency}
%%%%%%%%%%%%%%%%%%%%%%%%%%%%%%%%%%%%%%%%%%%%%%%%%%%%%%%%%%%%

The frequency of a value or category is the number of observations
having that value.

For the data

\[
4,\ 7,\ 6,\ 4,\ 9,\ 7,\ 5,\ 4,
\]

the value \(4\) occurs three times.

Thus,

\[
\boxed{
f(4)=3.
}
\]

%%%%%%%%%%%%%%%%%%%%%%%%%%%%%%%%%%%%%%%%%%%%%%%%%%%%%%%%%%%%
\subsection{Relative Frequency}
\label{subsec:relative-frequency}
%%%%%%%%%%%%%%%%%%%%%%%%%%%%%%%%%%%%%%%%%%%%%%%%%%%%%%%%%%%%

If category \(i\) occurs \(f_i\) times in a sample of size \(n\), its
relative frequency is

\[
\boxed{
r_i
=
\frac{f_i}{n}.
}
\]

Relative frequencies satisfy

\[
\boxed{
\sum_i r_i=1.
}
\]

Multiplying by \(100\%\) gives percentages.

%%%%%%%%%%%%%%%%%%%%%%%%%%%%%%%%%%%%%%%%%%%%%%%%%%%%%%%%%%%%
\subsection{Frequency Table}
\label{subsec:frequency-table}
%%%%%%%%%%%%%%%%%%%%%%%%%%%%%%%%%%%%%%%%%%%%%%%%%%%%%%%%%%%%

For the sample

\[
4,\ 7,\ 6,\ 4,\ 9,\ 7,\ 5,\ 4,
\]

a frequency table is

\[
\begin{array}{c|c|c}
x & f & f/8\\
\hline
4&3&3/8\\
5&1&1/8\\
6&1&1/8\\
7&2&2/8\\
9&1&1/8
\end{array}
\]

Frequency tables summarize repeated values efficiently.

%%%%%%%%%%%%%%%%%%%%%%%%%%%%%%%%%%%%%%%%%%%%%%%%%%%%%%%%%%%%
\subsection{Cumulative Frequency}
\label{subsec:cumulative-frequency}
%%%%%%%%%%%%%%%%%%%%%%%%%%%%%%%%%%%%%%%%%%%%%%%%%%%%%%%%%%%%

For ordered data, the cumulative frequency at \(x\) is the number of
observations less than or equal to \(x\).

One may define

\[
\boxed{
F_n(x)
=
\#\{i:x_i\leq x\}.
}
\]

Dividing by \(n\) gives cumulative relative frequency.

This leads naturally to the empirical distribution function.

%%%%%%%%%%%%%%%%%%%%%%%%%%%%%%%%%%%%%%%%%%%%%%%%%%%%%%%%%%%%
\subsection{Bar Charts}
\label{subsec:bar-charts}
%%%%%%%%%%%%%%%%%%%%%%%%%%%%%%%%%%%%%%%%%%%%%%%%%%%%%%%%%%%%

A bar chart displays categorical frequencies or relative frequencies.

Each bar corresponds to a category.

Bars are typically separated because the categories are distinct rather
than intervals on a continuous numerical scale.

Bar charts are especially useful for nominal and ordinal data.

%%%%%%%%%%%%%%%%%%%%%%%%%%%%%%%%%%%%%%%%%%%%%%%%%%%%%%%%%%%%
\subsection{Histograms}
\label{subsec:histograms}
%%%%%%%%%%%%%%%%%%%%%%%%%%%%%%%%%%%%%%%%%%%%%%%%%%%%%%%%%%%%

A histogram divides numerical data into intervals called bins.

The height or area of each bar represents the number or proportion of
observations in that interval.

Unlike a bar chart, adjacent histogram bars usually touch because the
horizontal axis represents a continuous numerical scale.

%%%%%%%%%%%%%%%%%%%%%%%%%%%%%%%%%%%%%%%%%%%%%%%%%%%%%%%%%%%%
\subsection{Histogram Bin Width}
\label{subsec:histogram-bin-width}
%%%%%%%%%%%%%%%%%%%%%%%%%%%%%%%%%%%%%%%%%%%%%%%%%%%%%%%%%%%%

The appearance of a histogram depends strongly on bin width.

Bins that are too wide may hide important structure.

Bins that are too narrow may create excessive noise.

Thus a histogram is not uniquely determined by the data.

It is a visualization whose construction involves a smoothing choice.

%%%%%%%%%%%%%%%%%%%%%%%%%%%%%%%%%%%%%%%%%%%%%%%%%%%%%%%%%%%%
\subsection{Density Histograms}
\label{subsec:density-histograms}
%%%%%%%%%%%%%%%%%%%%%%%%%%%%%%%%%%%%%%%%%%%%%%%%%%%%%%%%%%%%

A histogram may be scaled so that total bar area equals \(1\).

If a bin has width \(w_i\) and relative frequency \(r_i\), define its
density height as

\[
\boxed{
h_i
=
\frac{r_i}{w_i}.
}
\]

Then

\[
h_iw_i=r_i.
\]

The total area becomes

\[
\boxed{
\sum_i h_iw_i=1.
}
\]

Density histograms can be compared directly with probability densities.

%%%%%%%%%%%%%%%%%%%%%%%%%%%%%%%%%%%%%%%%%%%%%%%%%%%%%%%%%%%%
\subsection{Dotplots and Stem-and-Leaf Displays}
\label{subsec:dotplots-stem-leaf}
%%%%%%%%%%%%%%%%%%%%%%%%%%%%%%%%%%%%%%%%%%%%%%%%%%%%%%%%%%%%

For small or moderate datasets, displays that preserve individual
observations can be useful.

A dotplot places one mark for each observation.

A stem-and-leaf display separates each value into a leading stem and
trailing leaf.

These preserve more raw-data information than a histogram.

%%%%%%%%%%%%%%%%%%%%%%%%%%%%%%%%%%%%%%%%%%%%%%%%%%%%%%%%%%%%
\subsection{Measures of Center}
\label{subsec:measures-of-center}
%%%%%%%%%%%%%%%%%%%%%%%%%%%%%%%%%%%%%%%%%%%%%%%%%%%%%%%%%%%%

Measures of center describe a representative or central location of the
data.

Important examples are:

\[
\boxed{
\text{mean},
}
\]

\[
\boxed{
\text{median},
}
\]

and

\[
\boxed{
\text{mode}.
}
\]

No single measure is best for every dataset.

%%%%%%%%%%%%%%%%%%%%%%%%%%%%%%%%%%%%%%%%%%%%%%%%%%%%%%%%%%%%
\subsection{Arithmetic Mean}
\label{subsec:arithmetic-mean}
%%%%%%%%%%%%%%%%%%%%%%%%%%%%%%%%%%%%%%%%%%%%%%%%%%%%%%%%%%%%

\begin{definitionbox}{Sample Mean}
For observations

\[
x_1,\ldots,x_n,
\]

the arithmetic sample mean is

\[
\boxed{
\overline{x}
=
\frac1n
\sum_{i=1}^{n}x_i.
}
\]
\end{definitionbox}

The symbol

\[
\overline{x}
\]

is read ``x-bar.''

%%%%%%%%%%%%%%%%%%%%%%%%%%%%%%%%%%%%%%%%%%%%%%%%%%%%%%%%%%%%
\subsection{Population Mean}
\label{subsec:population-mean}
%%%%%%%%%%%%%%%%%%%%%%%%%%%%%%%%%%%%%%%%%%%%%%%%%%%%%%%%%%%%

For a finite population

\[
x_1,\ldots,x_N,
\]

the population mean is

\[
\boxed{
\mu
=
\frac1N
\sum_{i=1}^{N}x_i.
}
\]

The Greek letter

\[
\mu
\]

is called ``mu.''

%%%%%%%%%%%%%%%%%%%%%%%%%%%%%%%%%%%%%%%%%%%%%%%%%%%%%%%%%%%%
\subsection{Worked Example: Sample Mean}
\label{subsec:worked-sample-mean}
%%%%%%%%%%%%%%%%%%%%%%%%%%%%%%%%%%%%%%%%%%%%%%%%%%%%%%%%%%%%

\begin{workedexamplebox}{Computing a Mean}

Consider

\[
4,\ 7,\ 6,\ 4,\ 9,\ 7,\ 5,\ 4.
\]

Then

\[
n=8.
\]

The sum is

\[
4+7+6+4+9+7+5+4=46.
\]

Therefore,

\begin{answerbox}
\[
\boxed{
\overline{x}
=
\frac{46}{8}
=
5.75.
}
\]
\end{answerbox}

\end{workedexamplebox}

%%%%%%%%%%%%%%%%%%%%%%%%%%%%%%%%%%%%%%%%%%%%%%%%%%%%%%%%%%%%
\subsection{Mean as a Balance Point}
\label{subsec:mean-balance-point}
%%%%%%%%%%%%%%%%%%%%%%%%%%%%%%%%%%%%%%%%%%%%%%%%%%%%%%%%%%%%

The mean satisfies

\[
\boxed{
\sum_{i=1}^{n}
(x_i-\overline{x})
=
0.
}
\]

Thus positive deviations above the mean balance negative deviations below
the mean.

This gives a geometric interpretation of the mean as a balance point.

%%%%%%%%%%%%%%%%%%%%%%%%%%%%%%%%%%%%%%%%%%%%%%%%%%%%%%%%%%%%
\subsection{Weighted Mean}
\label{subsec:weighted-mean}
%%%%%%%%%%%%%%%%%%%%%%%%%%%%%%%%%%%%%%%%%%%%%%%%%%%%%%%%%%%%

Suppose values \(x_i\) have nonnegative weights \(w_i\).

The weighted mean is

\[
\boxed{
\overline{x}_w
=
\frac{
\sum_iw_ix_i
}{
\sum_iw_i
}.
}
\]

If

\[
\sum_iw_i=1,
\]

then

\[
\boxed{
\overline{x}_w
=
\sum_iw_ix_i.
}
\]

%%%%%%%%%%%%%%%%%%%%%%%%%%%%%%%%%%%%%%%%%%%%%%%%%%%%%%%%%%%%
\subsection{Worked Example: Weighted Mean}
\label{subsec:worked-weighted-mean}
%%%%%%%%%%%%%%%%%%%%%%%%%%%%%%%%%%%%%%%%%%%%%%%%%%%%%%%%%%%%

\begin{workedexamplebox}{Weighted Course Average}

Suppose exam scores are

\[
80,\ 90,\ 70
\]

with weights

\[
0.30,\ 0.40,\ 0.30.
\]

Then

\[
\overline{x}_w
=
0.30(80)
+
0.40(90)
+
0.30(70).
\]

Thus,

\begin{answerbox}
\[
\boxed{
\overline{x}_w
=
81.
}
\]
\end{answerbox}

\end{workedexamplebox}

%%%%%%%%%%%%%%%%%%%%%%%%%%%%%%%%%%%%%%%%%%%%%%%%%%%%%%%%%%%%
\subsection{Median}
\label{subsec:median}
%%%%%%%%%%%%%%%%%%%%%%%%%%%%%%%%%%%%%%%%%%%%%%%%%%%%%%%%%%%%

The median is the middle ordered observation.

Sort the data:

\[
x_{(1)}
\leq
x_{(2)}
\leq
\cdots
\leq
x_{(n)}.
\]

If \(n\) is odd,

\[
\boxed{
\operatorname{Median}
=
x_{\left(\frac{n+1}{2}\right)}.
}
\]

If \(n\) is even, a common convention is

\[
\boxed{
\operatorname{Median}
=
\frac{
x_{(n/2)}
+
x_{(n/2+1)}
}{2}.
}
\]

%%%%%%%%%%%%%%%%%%%%%%%%%%%%%%%%%%%%%%%%%%%%%%%%%%%%%%%%%%%%
\subsection{Worked Example: Median}
\label{subsec:worked-median}
%%%%%%%%%%%%%%%%%%%%%%%%%%%%%%%%%%%%%%%%%%%%%%%%%%%%%%%%%%%%

\begin{workedexamplebox}{Finding the Median}

Sort the sample

\[
4,\ 7,\ 6,\ 4,\ 9,\ 7,\ 5,\ 4
\]

to obtain

\[
4,\ 4,\ 4,\ 5,\ 6,\ 7,\ 7,\ 9.
\]

Since \(n=8\), average the fourth and fifth observations:

\[
\frac{5+6}{2}.
\]

Therefore,

\begin{answerbox}
\[
\boxed{
\operatorname{Median}
=
5.5.
}
\]
\end{answerbox}

\end{workedexamplebox}

%%%%%%%%%%%%%%%%%%%%%%%%%%%%%%%%%%%%%%%%%%%%%%%%%%%%%%%%%%%%
\subsection{Mode}
\label{subsec:mode}
%%%%%%%%%%%%%%%%%%%%%%%%%%%%%%%%%%%%%%%%%%%%%%%%%%%%%%%%%%%%

The mode is the most frequently occurring value or category.

A dataset may have:

\[
\boxed{
\text{one mode},
}
\]

\[
\boxed{
\text{multiple modes},
}
\]

or

\[
\boxed{
\text{no unique mode}.
}
\]

For

\[
4,\ 4,\ 4,\ 5,\ 6,\ 7,\ 7,\ 9,
\]

the mode is

\[
\boxed{
4.
}
\]

%%%%%%%%%%%%%%%%%%%%%%%%%%%%%%%%%%%%%%%%%%%%%%%%%%%%%%%%%%%%
\subsection{Mean Versus Median}
\label{subsec:mean-versus-median}
%%%%%%%%%%%%%%%%%%%%%%%%%%%%%%%%%%%%%%%%%%%%%%%%%%%%%%%%%%%%

The mean uses every numerical value.

Therefore it is sensitive to extreme observations.

The median depends mainly on the ordering of observations.

Therefore it is more resistant to extreme values.

For strongly skewed data,

\[
\boxed{
\text{median}
}
\]

may better represent a typical observation.

%%%%%%%%%%%%%%%%%%%%%%%%%%%%%%%%%%%%%%%%%%%%%%%%%%%%%%%%%%%%
\subsection{Worked Example: Effect of an Outlier}
\label{subsec:worked-outlier-center}
%%%%%%%%%%%%%%%%%%%%%%%%%%%%%%%%%%%%%%%%%%%%%%%%%%%%%%%%%%%%

\begin{workedexamplebox}{Mean and Median Under an Extreme Value}

Consider

\[
2,\ 3,\ 4,\ 5,\ 6.
\]

Both mean and median are

\[
4.
\]

Replace \(6\) by \(100\):

\[
2,\ 3,\ 4,\ 5,\ 100.
\]

The median remains

\[
4.
\]

The mean becomes

\[
\frac{114}{5}=22.8.
\]

\begin{answerbox}
\[
\boxed{
\text{The median is much more resistant to the outlier.}
}
\]
\end{answerbox}

\end{workedexamplebox}

%%%%%%%%%%%%%%%%%%%%%%%%%%%%%%%%%%%%%%%%%%%%%%%%%%%%%%%%%%%%
\subsection{Measures of Spread}
\label{subsec:measures-of-spread}
%%%%%%%%%%%%%%%%%%%%%%%%%%%%%%%%%%%%%%%%%%%%%%%%%%%%%%%%%%%%

Measures of spread describe how dispersed observations are.

Important measures include:

\[
\boxed{
\text{range},
}
\]

\[
\boxed{
\text{interquartile range},
}
\]

\[
\boxed{
\text{variance},
}
\]

and

\[
\boxed{
\text{standard deviation}.
}
\]

%%%%%%%%%%%%%%%%%%%%%%%%%%%%%%%%%%%%%%%%%%%%%%%%%%%%%%%%%%%%
\subsection{Range}
\label{subsec:range}
%%%%%%%%%%%%%%%%%%%%%%%%%%%%%%%%%%%%%%%%%%%%%%%%%%%%%%%%%%%%

The range is

\[
\boxed{
R
=
x_{\max}-x_{\min}.
}
\]

It uses only the smallest and largest observations.

Therefore it is highly sensitive to extreme values.

%%%%%%%%%%%%%%%%%%%%%%%%%%%%%%%%%%%%%%%%%%%%%%%%%%%%%%%%%%%%
\subsection{Deviations from the Mean}
\label{subsec:deviations-from-mean}
%%%%%%%%%%%%%%%%%%%%%%%%%%%%%%%%%%%%%%%%%%%%%%%%%%%%%%%%%%%%

The deviation of observation \(x_i\) from the sample mean is

\[
\boxed{
d_i
=
x_i-\overline{x}.
}
\]

Because

\[
\sum_i
(x_i-\overline{x})
=
0,
\]

simple averaging of signed deviations cannot measure spread.

One therefore uses absolute or squared deviations.

%%%%%%%%%%%%%%%%%%%%%%%%%%%%%%%%%%%%%%%%%%%%%%%%%%%%%%%%%%%%
\subsection{Population Variance}
\label{subsec:population-variance-statistics}
%%%%%%%%%%%%%%%%%%%%%%%%%%%%%%%%%%%%%%%%%%%%%%%%%%%%%%%%%%%%

For a finite population,

\[
\boxed{
\sigma^2
=
\frac1N
\sum_{i=1}^{N}
(x_i-\mu)^2.
}
\]

The population standard deviation is

\[
\boxed{
\sigma
=
\sqrt{\sigma^2}.
}
\]

Variance measures average squared distance from the population mean.

%%%%%%%%%%%%%%%%%%%%%%%%%%%%%%%%%%%%%%%%%%%%%%%%%%%%%%%%%%%%
\subsection{Sample Variance}
\label{subsec:sample-variance}
%%%%%%%%%%%%%%%%%%%%%%%%%%%%%%%%%%%%%%%%%%%%%%%%%%%%%%%%%%%%

\begin{definitionbox}{Sample Variance}
For observations

\[
x_1,\ldots,x_n,
\]

the sample variance is

\[
\boxed{
s^2
=
\frac1{n-1}
\sum_{i=1}^{n}
(x_i-\overline{x})^2.
}
\]
\end{definitionbox}

The sample standard deviation is

\[
\boxed{
s
=
\sqrt{s^2}.
}
\]

%%%%%%%%%%%%%%%%%%%%%%%%%%%%%%%%%%%%%%%%%%%%%%%%%%%%%%%%%%%%
\subsection{Why Divide by \(n-1\)?}
\label{subsec:why-n-minus-one}
%%%%%%%%%%%%%%%%%%%%%%%%%%%%%%%%%%%%%%%%%%%%%%%%%%%%%%%%%%%%

Once the sample mean is known, the deviations satisfy

\[
\sum_{i=1}^{n}
(x_i-\overline{x})
=
0.
\]

Therefore only \(n-1\) deviations can vary freely.

This motivates the phrase

\[
\boxed{
n-1
\text{ degrees of freedom}.
}
\]

More importantly, when observations are i.i.d. with population variance
\(\sigma^2\),

\[
\boxed{
\mathbb{E}[s^2]
=
\sigma^2.
}
\]

Thus division by \(n-1\) makes the usual sample variance unbiased for
the population variance.

%%%%%%%%%%%%%%%%%%%%%%%%%%%%%%%%%%%%%%%%%%%%%%%%%%%%%%%%%%%%
\subsection{Computational Formula for Variance}
\label{subsec:computational-variance-formula}
%%%%%%%%%%%%%%%%%%%%%%%%%%%%%%%%%%%%%%%%%%%%%%%%%%%%%%%%%%%%

Using

\[
\sum_{i=1}^{n}
(x_i-\overline{x})^2
=
\sum_{i=1}^{n}x_i^2
-
n\overline{x}^2,
\]

we obtain

\[
\boxed{
s^2
=
\frac{
\sum x_i^2
-
n\overline{x}^2
}{
n-1
}.
}
\]

Equivalently,

\[
\boxed{
s^2
=
\frac{
\sum x_i^2
-
(\sum x_i)^2/n
}{
n-1
}.
}
\]

%%%%%%%%%%%%%%%%%%%%%%%%%%%%%%%%%%%%%%%%%%%%%%%%%%%%%%%%%%%%
\subsection{Worked Example: Sample Variance}
\label{subsec:worked-sample-variance}
%%%%%%%%%%%%%%%%%%%%%%%%%%%%%%%%%%%%%%%%%%%%%%%%%%%%%%%%%%%%

\begin{workedexamplebox}{Computing Variance and Standard Deviation}

Consider

\[
2,\ 4,\ 6.
\]

The sample mean is

\[
\overline{x}=4.
\]

The squared deviations are

\[
(2-4)^2=4,
\]

\[
(4-4)^2=0,
\]

and

\[
(6-4)^2=4.
\]

Thus,

\[
\sum
(x_i-\overline{x})^2
=
8.
\]

The sample variance is

\[
s^2
=
\frac8{3-1}
=
4.
\]

Therefore,

\begin{answerbox}
\[
\boxed{
s^2=4,
\qquad
s=2.
}
\]
\end{answerbox}

\end{workedexamplebox}

%%%%%%%%%%%%%%%%%%%%%%%%%%%%%%%%%%%%%%%%%%%%%%%%%%%%%%%%%%%%
\subsection{Units of Variance and Standard Deviation}
\label{subsec:variance-units}
%%%%%%%%%%%%%%%%%%%%%%%%%%%%%%%%%%%%%%%%%%%%%%%%%%%%%%%%%%%%

If observations are measured in meters, then variance is measured in

\[
\text{meters}^2.
\]

Standard deviation returns to the original units:

\[
\boxed{
\text{SD has the same units as the data}.
}
\]

This is one reason standard deviation is often easier to interpret than
variance.

%%%%%%%%%%%%%%%%%%%%%%%%%%%%%%%%%%%%%%%%%%%%%%%%%%%%%%%%%%%%
\subsection{Interquartile Range}
\label{subsec:interquartile-range}
%%%%%%%%%%%%%%%%%%%%%%%%%%%%%%%%%%%%%%%%%%%%%%%%%%%%%%%%%%%%

Let

\[
Q_1
\]

be the first quartile and

\[
Q_3
\]

the third quartile.

The interquartile range is

\[
\boxed{
IQR
=
Q_3-Q_1.
}
\]

It measures the spread of the middle \(50\%\) of the data.

Because it ignores the extreme tails, it is more resistant than the
range or standard deviation.

%%%%%%%%%%%%%%%%%%%%%%%%%%%%%%%%%%%%%%%%%%%%%%%%%%%%%%%%%%%%
\subsection{Quartiles}
\label{subsec:quartiles}
%%%%%%%%%%%%%%%%%%%%%%%%%%%%%%%%%%%%%%%%%%%%%%%%%%%%%%%%%%%%

Quartiles divide ordered data into approximately four equal portions.

The common notation is:

\[
\boxed{
Q_1
=
25\text{th percentile},
}
\]

\[
\boxed{
Q_2
=
\text{median},
}
\]

and

\[
\boxed{
Q_3
=
75\text{th percentile}.
}
\]

Exact quartile interpolation conventions vary among textbooks and
software packages.

%%%%%%%%%%%%%%%%%%%%%%%%%%%%%%%%%%%%%%%%%%%%%%%%%%%%%%%%%%%%
\subsection{Percentiles}
\label{subsec:percentiles}
%%%%%%%%%%%%%%%%%%%%%%%%%%%%%%%%%%%%%%%%%%%%%%%%%%%%%%%%%%%%

The \(p\)-th percentile is a value below which approximately \(p\%\) of
the observations fall.

Thus the \(90\)-th percentile represents a relatively high observation.

Percentiles describe position relative to the entire sample.

%%%%%%%%%%%%%%%%%%%%%%%%%%%%%%%%%%%%%%%%%%%%%%%%%%%%%%%%%%%%
\subsection{Quantiles}
\label{subsec:sample-quantiles}
%%%%%%%%%%%%%%%%%%%%%%%%%%%%%%%%%%%%%%%%%%%%%%%%%%%%%%%%%%%%

For

\[
0\leq q\leq1,
\]

a \(q\)-quantile is a value corresponding to cumulative proportion \(q\).

Examples include:

\[
q=0.25
\longrightarrow
Q_1,
\]

\[
q=0.50
\longrightarrow
\text{median},
\]

and

\[
q=0.75
\longrightarrow
Q_3.
\]

Sample quantile definitions may differ slightly because finite datasets
do not always divide evenly.

%%%%%%%%%%%%%%%%%%%%%%%%%%%%%%%%%%%%%%%%%%%%%%%%%%%%%%%%%%%%
\subsection{Five-Number Summary}
\label{subsec:five-number-summary}
%%%%%%%%%%%%%%%%%%%%%%%%%%%%%%%%%%%%%%%%%%%%%%%%%%%%%%%%%%%%

The five-number summary consists of

\[
\boxed{
\text{minimum},
}
\]

\[
\boxed{
Q_1,
}
\]

\[
\boxed{
\text{median},
}
\]

\[
\boxed{
Q_3,
}
\]

and

\[
\boxed{
\text{maximum}.
}
\]

It provides a compact resistant summary of location and spread.

%%%%%%%%%%%%%%%%%%%%%%%%%%%%%%%%%%%%%%%%%%%%%%%%%%%%%%%%%%%%
\subsection{Outlier Fences}
\label{subsec:outlier-fences}
%%%%%%%%%%%%%%%%%%%%%%%%%%%%%%%%%%%%%%%%%%%%%%%%%%%%%%%%%%%%

A common exploratory rule defines the lower fence

\[
\boxed{
Q_1-1.5(IQR)
}
\]

and upper fence

\[
\boxed{
Q_3+1.5(IQR).
}
\]

Observations outside these fences are flagged as potential outliers.

This is an exploratory rule, not a proof that an observation is
erroneous.

%%%%%%%%%%%%%%%%%%%%%%%%%%%%%%%%%%%%%%%%%%%%%%%%%%%%%%%%%%%%
\subsection{Extreme Outlier Fences}
\label{subsec:extreme-outlier-fences}
%%%%%%%%%%%%%%%%%%%%%%%%%%%%%%%%%%%%%%%%%%%%%%%%%%%%%%%%%%%%

A more conservative threshold sometimes uses

\[
\boxed{
Q_1-3(IQR)
}
\]

and

\[
\boxed{
Q_3+3(IQR).
}
\]

Observations outside these limits may be described as extreme outliers.

Again, context determines whether an unusual observation should be
removed, corrected, or retained.

%%%%%%%%%%%%%%%%%%%%%%%%%%%%%%%%%%%%%%%%%%%%%%%%%%%%%%%%%%%%
\subsection{Boxplots}
\label{subsec:boxplots}
%%%%%%%%%%%%%%%%%%%%%%%%%%%%%%%%%%%%%%%%%%%%%%%%%%%%%%%%%%%%

A boxplot graphically displays:

\[
Q_1,
\]

\[
\text{median},
\]

\[
Q_3,
\]

and whiskers extending to selected non-outlying observations.

Potential outliers are often displayed individually.

Boxplots are especially useful for comparing distributions across
several groups.

%%%%%%%%%%%%%%%%%%%%%%%%%%%%%%%%%%%%%%%%%%%%%%%%%%%%%%%%%%%%
\subsection{Worked Example: Outlier Rule}
\label{subsec:worked-outlier-rule}
%%%%%%%%%%%%%%%%%%%%%%%%%%%%%%%%%%%%%%%%%%%%%%%%%%%%%%%%%%%%

\begin{workedexamplebox}{Using the \(1.5IQR\) Rule}

Suppose

\[
Q_1=10
\]

and

\[
Q_3=18.
\]

Then

\[
IQR=8.
\]

The lower fence is

\[
10-1.5(8)
=
-2.
\]

The upper fence is

\[
18+1.5(8)
=
30.
\]

Therefore,

\begin{answerbox}
\[
\boxed{
\text{values below }-2
\text{ or above }30
\text{ are flagged as possible outliers}.
}
\]
\end{answerbox}

\end{workedexamplebox}

%%%%%%%%%%%%%%%%%%%%%%%%%%%%%%%%%%%%%%%%%%%%%%%%%%%%%%%%%%%%
\subsection{Standardized Scores}
\label{subsec:standardized-scores}
%%%%%%%%%%%%%%%%%%%%%%%%%%%%%%%%%%%%%%%%%%%%%%%%%%%%%%%%%%%%

A standardized score describes an observation relative to a mean and
standard deviation.

For population quantities,

\[
\boxed{
z
=
\frac{x-\mu}{\sigma}.
}
\]

For sample quantities,

\[
\boxed{
z
=
\frac{x-\overline{x}}{s}.
}
\]

The standardized value is called a \(z\)-score.

%%%%%%%%%%%%%%%%%%%%%%%%%%%%%%%%%%%%%%%%%%%%%%%%%%%%%%%%%%%%
\subsection{Interpretation of a \(z\)-Score}
\label{subsec:z-score-interpretation}
%%%%%%%%%%%%%%%%%%%%%%%%%%%%%%%%%%%%%%%%%%%%%%%%%%%%%%%%%%%%

If

\[
z=2,
\]

the observation lies two standard deviations above the mean.

If

\[
z=-1.5,
\]

the observation lies \(1.5\) standard deviations below the mean.

A \(z\)-score has no physical units.

This allows comparisons across differently scaled variables.

%%%%%%%%%%%%%%%%%%%%%%%%%%%%%%%%%%%%%%%%%%%%%%%%%%%%%%%%%%%%
\subsection{Worked Example: \(z\)-Score}
\label{subsec:worked-z-score}
%%%%%%%%%%%%%%%%%%%%%%%%%%%%%%%%%%%%%%%%%%%%%%%%%%%%%%%%%%%%

\begin{workedexamplebox}{Standardizing an Exam Score}

Suppose

\[
\mu=70,
\qquad
\sigma=8,
\]

and a student scores

\[
x=86.
\]

Then

\[
z
=
\frac{86-70}{8}.
\]

Therefore,

\begin{answerbox}
\[
\boxed{
z=2.
}
\]
\end{answerbox}

The score is two standard deviations above the mean.

\end{workedexamplebox}

%%%%%%%%%%%%%%%%%%%%%%%%%%%%%%%%%%%%%%%%%%%%%%%%%%%%%%%%%%%%
\subsection{Effect of Linear Transformations}
\label{subsec:linear-transformations-descriptive}
%%%%%%%%%%%%%%%%%%%%%%%%%%%%%%%%%%%%%%%%%%%%%%%%%%%%%%%%%%%%

Suppose all observations are transformed by

\[
\boxed{
y_i=ax_i+b.
}
\]

Then the mean transforms as

\[
\boxed{
\overline{y}
=
a\overline{x}+b.
}
\]

The variance transforms as

\[
\boxed{
s_y^2
=
a^2s_x^2.
}
\]

The standard deviation transforms as

\[
\boxed{
s_y
=
|a|s_x.
}
\]

%%%%%%%%%%%%%%%%%%%%%%%%%%%%%%%%%%%%%%%%%%%%%%%%%%%%%%%%%%%%
\subsection{Effect on Median and Quantiles}
\label{subsec:transform-median-quantiles}
%%%%%%%%%%%%%%%%%%%%%%%%%%%%%%%%%%%%%%%%%%%%%%%%%%%%%%%%%%%%

For \(a>0\),

\[
y=ax+b
\]

preserves ordering.

Thus quantiles transform as

\[
\boxed{
Q_y(q)
=
aQ_x(q)+b.
}
\]

For \(a<0\), ordering reverses, so lower and upper quantiles exchange
roles appropriately.

%%%%%%%%%%%%%%%%%%%%%%%%%%%%%%%%%%%%%%%%%%%%%%%%%%%%%%%%%%%%
\subsection{Standardization Properties}
\label{subsec:standardization-properties}
%%%%%%%%%%%%%%%%%%%%%%%%%%%%%%%%%%%%%%%%%%%%%%%%%%%%%%%%%%%%

If

\[
z_i
=
\frac{x_i-\overline{x}}{s},
\]

then the standardized sample has mean

\[
\boxed{
\overline{z}=0.
}
\]

Using the same \(n-1\) sample-variance convention, its sample standard
deviation is

\[
\boxed{
s_z=1.
}
\]

Thus standardization removes location and scale.

%%%%%%%%%%%%%%%%%%%%%%%%%%%%%%%%%%%%%%%%%%%%%%%%%%%%%%%%%%%%
\subsection{Shape of a Distribution}
\label{subsec:shape-distribution-descriptive}
%%%%%%%%%%%%%%%%%%%%%%%%%%%%%%%%%%%%%%%%%%%%%%%%%%%%%%%%%%%%

Beyond center and spread, important shape characteristics include:

\[
\boxed{
\text{symmetry},
}
\]

\[
\boxed{
\text{skewness},
}
\]

\[
\boxed{
\text{modality},
}
\]

and

\[
\boxed{
\text{tail behavior}.
}
\]

Graphical displays are often necessary for understanding shape.

%%%%%%%%%%%%%%%%%%%%%%%%%%%%%%%%%%%%%%%%%%%%%%%%%%%%%%%%%%%%
\subsection{Symmetric Data}
\label{subsec:symmetric-data}
%%%%%%%%%%%%%%%%%%%%%%%%%%%%%%%%%%%%%%%%%%%%%%%%%%%%%%%%%%%%

A distribution is symmetric around \(c\) if its left and right sides
mirror one another.

For an exactly symmetric unimodal distribution, the mean and median are
often equal:

\[
\boxed{
\text{mean}
=
\text{median}.
}
\]

The normal distribution is the most familiar symmetric example.

%%%%%%%%%%%%%%%%%%%%%%%%%%%%%%%%%%%%%%%%%%%%%%%%%%%%%%%%%%%%
\subsection{Right Skew}
\label{subsec:right-skew}
%%%%%%%%%%%%%%%%%%%%%%%%%%%%%%%%%%%%%%%%%%%%%%%%%%%%%%%%%%%%

A right-skewed distribution has a long tail extending toward larger
values.

Such data may arise from:

\[
\text{income},
\]

\[
\text{waiting times},
\]

or

\[
\text{repair costs}.
\]

Large values tend to pull the mean to the right.

Therefore one often sees

\[
\boxed{
\text{mean}
>
\text{median}.
}
\]

This is a tendency rather than an absolute mathematical rule for every
finite dataset.

%%%%%%%%%%%%%%%%%%%%%%%%%%%%%%%%%%%%%%%%%%%%%%%%%%%%%%%%%%%%
\subsection{Left Skew}
\label{subsec:left-skew}
%%%%%%%%%%%%%%%%%%%%%%%%%%%%%%%%%%%%%%%%%%%%%%%%%%%%%%%%%%%%

A left-skewed distribution has a long tail toward smaller values.

Small extreme observations pull the mean downward.

Thus one often sees

\[
\boxed{
\text{mean}
<
\text{median}.
}
\]

Again, this is a descriptive tendency rather than a universal identity.

%%%%%%%%%%%%%%%%%%%%%%%%%%%%%%%%%%%%%%%%%%%%%%%%%%%%%%%%%%%%
\subsection{Sample Moments}
\label{subsec:sample-moments}
%%%%%%%%%%%%%%%%%%%%%%%%%%%%%%%%%%%%%%%%%%%%%%%%%%%%%%%%%%%%

The \(k\)-th sample raw moment may be defined as

\[
\boxed{
m_k'
=
\frac1n
\sum_{i=1}^{n}
x_i^k.
}
\]

The \(k\)-th sample central moment is

\[
\boxed{
m_k
=
\frac1n
\sum_{i=1}^{n}
(x_i-\overline{x})^k.
}
\]

These empirical moments parallel the moments of probability
distributions.

%%%%%%%%%%%%%%%%%%%%%%%%%%%%%%%%%%%%%%%%%%%%%%%%%%%%%%%%%%%%
\subsection{Sample Skewness}
\label{subsec:sample-skewness}
%%%%%%%%%%%%%%%%%%%%%%%%%%%%%%%%%%%%%%%%%%%%%%%%%%%%%%%%%%%%

A moment-based measure of skewness uses the standardized third central
moment.

One simple descriptive version is

\[
\boxed{
g_1
=
\frac{
\frac1n\sum_{i=1}^{n}
(x_i-\overline{x})^3
}{
\left[
\frac1n\sum_{i=1}^{n}
(x_i-\overline{x})^2
\right]^{3/2}
}.
}
\]

Positive values indicate right-skew tendency.

Negative values indicate left-skew tendency.

Different statistical software may use bias-corrected versions.

%%%%%%%%%%%%%%%%%%%%%%%%%%%%%%%%%%%%%%%%%%%%%%%%%%%%%%%%%%%%
\subsection{Kurtosis}
\label{subsec:sample-kurtosis}
%%%%%%%%%%%%%%%%%%%%%%%%%%%%%%%%%%%%%%%%%%%%%%%%%%%%%%%%%%%%

A standardized fourth-moment measure is

\[
\boxed{
g_2
=
\frac{
\frac1n\sum_{i=1}^{n}
(x_i-\overline{x})^4
}{
\left[
\frac1n\sum_{i=1}^{n}
(x_i-\overline{x})^2
\right]^2
}.
}
\]

For a normal distribution, the population kurtosis is

\[
3.
\]

Excess kurtosis subtracts \(3\):

\[
\boxed{
\text{excess kurtosis}
=
g_2-3
}
\]

under the corresponding moment convention.

%%%%%%%%%%%%%%%%%%%%%%%%%%%%%%%%%%%%%%%%%%%%%%%%%%%%%%%%%%%%
\subsection{Interpreting Kurtosis Carefully}
\label{subsec:kurtosis-interpretation}
%%%%%%%%%%%%%%%%%%%%%%%%%%%%%%%%%%%%%%%%%%%%%%%%%%%%%%%%%%%%

Kurtosis is often described informally as ``peakedness,'' but this can
be misleading.

The fourth moment is highly sensitive to large deviations.

Thus kurtosis is better understood as reflecting the influence of the
tails and extreme observations in a standardized fourth-moment sense.

%%%%%%%%%%%%%%%%%%%%%%%%%%%%%%%%%%%%%%%%%%%%%%%%%%%%%%%%%%%%
\subsection{Empirical Distribution Function}
\label{subsec:empirical-distribution-function}
%%%%%%%%%%%%%%%%%%%%%%%%%%%%%%%%%%%%%%%%%%%%%%%%%%%%%%%%%%%%

\begin{definitionbox}{Empirical Distribution Function}
For data

\[
x_1,\ldots,x_n,
\]

the empirical distribution function is

\[
\boxed{
F_n(x)
=
\frac1n
\sum_{i=1}^{n}
\mathbf{1}_{\{x_i\leq x\}}.
}
\]
\end{definitionbox}

Thus \(F_n(x)\) is the fraction of observations less than or equal to
\(x\).

%%%%%%%%%%%%%%%%%%%%%%%%%%%%%%%%%%%%%%%%%%%%%%%%%%%%%%%%%%%%
\subsection{Properties of the Empirical CDF}
\label{subsec:empirical-cdf-properties}
%%%%%%%%%%%%%%%%%%%%%%%%%%%%%%%%%%%%%%%%%%%%%%%%%%%%%%%%%%%%

The empirical CDF is:

\[
\boxed{
\text{nondecreasing},
}
\]

\[
\boxed{
\text{right-continuous},
}
\]

and takes values in

\[
[0,1].
\]

It jumps by multiples of

\[
1/n.
\]

At each observation, the cumulative proportion increases.

%%%%%%%%%%%%%%%%%%%%%%%%%%%%%%%%%%%%%%%%%%%%%%%%%%%%%%%%%%%%
\subsection{Worked Example: Empirical CDF}
\label{subsec:worked-empirical-cdf}
%%%%%%%%%%%%%%%%%%%%%%%%%%%%%%%%%%%%%%%%%%%%%%%%%%%%%%%%%%%%

\begin{workedexamplebox}{Evaluating an Empirical Distribution}

Suppose

\[
x_1,\ldots,x_5
=
2,\ 4,\ 4,\ 7,\ 10.
\]

Find

\[
F_5(4).
\]

Three observations satisfy

\[
x_i\leq4.
\]

Therefore,

\begin{answerbox}
\[
\boxed{
F_5(4)
=
\frac35.
}
\]
\end{answerbox}

\end{workedexamplebox}

%%%%%%%%%%%%%%%%%%%%%%%%%%%%%%%%%%%%%%%%%%%%%%%%%%%%%%%%%%%%
\subsection{Bivariate Data}
\label{subsec:bivariate-data}
%%%%%%%%%%%%%%%%%%%%%%%%%%%%%%%%%%%%%%%%%%%%%%%%%%%%%%%%%%%%

Suppose each observation consists of a pair

\[
(x_i,y_i).
\]

Then the data are bivariate.

Examples include:

\[
(\text{height},\text{weight}),
\]

\[
(\text{study time},\text{exam score}),
\]

or

\[
(\text{temperature},\text{energy use}).
\]

The goal is often to describe association between the two variables.

%%%%%%%%%%%%%%%%%%%%%%%%%%%%%%%%%%%%%%%%%%%%%%%%%%%%%%%%%%%%
\subsection{Scatterplots}
\label{subsec:scatterplots}
%%%%%%%%%%%%%%%%%%%%%%%%%%%%%%%%%%%%%%%%%%%%%%%%%%%%%%%%%%%%

A scatterplot places the point

\[
(x_i,y_i)
\]

for each observation.

It can reveal:

\[
\boxed{
\text{direction},
}
\]

\[
\boxed{
\text{strength},
}
\]

\[
\boxed{
\text{form},
}
\]

and

\[
\boxed{
\text{outliers}.
}
\]

Scatterplots should generally be inspected before computing correlation.

%%%%%%%%%%%%%%%%%%%%%%%%%%%%%%%%%%%%%%%%%%%%%%%%%%%%%%%%%%%%
\subsection{Sample Covariance}
\label{subsec:sample-covariance-statistics}
%%%%%%%%%%%%%%%%%%%%%%%%%%%%%%%%%%%%%%%%%%%%%%%%%%%%%%%%%%%%

For paired data

\[
(x_i,y_i),
\]

the sample covariance is

\[
\boxed{
s_{xy}
=
\frac1{n-1}
\sum_{i=1}^{n}
(x_i-\overline{x})
(y_i-\overline{y}).
}
\]

Positive covariance indicates a tendency for the variables to move
together.

Negative covariance indicates opposite movement.

%%%%%%%%%%%%%%%%%%%%%%%%%%%%%%%%%%%%%%%%%%%%%%%%%%%%%%%%%%%%
\subsection{Units of Covariance}
\label{subsec:covariance-units}
%%%%%%%%%%%%%%%%%%%%%%%%%%%%%%%%%%%%%%%%%%%%%%%%%%%%%%%%%%%%

Covariance depends on measurement units.

If \(X\) is measured in meters and \(Y\) in seconds, covariance has units

\[
\text{meter}\cdot\text{second}.
\]

Therefore the magnitude of covariance is difficult to compare across
differently scaled variables.

Correlation solves this by standardizing.

%%%%%%%%%%%%%%%%%%%%%%%%%%%%%%%%%%%%%%%%%%%%%%%%%%%%%%%%%%%%
\subsection{Sample Correlation}
\label{subsec:sample-correlation}
%%%%%%%%%%%%%%%%%%%%%%%%%%%%%%%%%%%%%%%%%%%%%%%%%%%%%%%%%%%%

\begin{definitionbox}{Sample Correlation}
The sample Pearson correlation coefficient is

\[
\boxed{
r
=
\frac{
s_{xy}
}{
s_xs_y
}.
}
\]
\end{definitionbox}

Equivalently,

\[
\boxed{
r
=
\frac{
\sum
(x_i-\overline{x})
(y_i-\overline{y})
}{
\sqrt{
\sum
(x_i-\overline{x})^2
}
\sqrt{
\sum
(y_i-\overline{y})^2
}
}.
}
\]

%%%%%%%%%%%%%%%%%%%%%%%%%%%%%%%%%%%%%%%%%%%%%%%%%%%%%%%%%%%%
\subsection{Correlation Bounds}
\label{subsec:correlation-bounds}
%%%%%%%%%%%%%%%%%%%%%%%%%%%%%%%%%%%%%%%%%%%%%%%%%%%%%%%%%%%%

The correlation coefficient satisfies

\[
\boxed{
-1\leq r\leq1.
}
\]

Values near \(1\) indicate strong positive linear association.

Values near \(-1\) indicate strong negative linear association.

Values near \(0\) indicate weak linear association.

%%%%%%%%%%%%%%%%%%%%%%%%%%%%%%%%%%%%%%%%%%%%%%%%%%%%%%%%%%%%
\subsection{Perfect Linear Correlation}
\label{subsec:perfect-linear-correlation}
%%%%%%%%%%%%%%%%%%%%%%%%%%%%%%%%%%%%%%%%%%%%%%%%%%%%%%%%%%%%

If

\[
r=1,
\]

all observations lie exactly on a line with positive slope.

If

\[
r=-1,
\]

all observations lie exactly on a line with negative slope.

Thus

\[
\boxed{
|r|=1
}
\]

corresponds to perfect linear association, assuming both variables have
nonzero sample standard deviations.

%%%%%%%%%%%%%%%%%%%%%%%%%%%%%%%%%%%%%%%%%%%%%%%%%%%%%%%%%%%%
\subsection{Correlation and Linear Transformations}
\label{subsec:correlation-transformations}
%%%%%%%%%%%%%%%%%%%%%%%%%%%%%%%%%%%%%%%%%%%%%%%%%%%%%%%%%%%%

Adding constants does not change correlation.

Multiplying either variable by a positive constant does not change
correlation.

Multiplying one variable by a negative constant reverses the sign.

Thus correlation measures scale-free linear association.

%%%%%%%%%%%%%%%%%%%%%%%%%%%%%%%%%%%%%%%%%%%%%%%%%%%%%%%%%%%%
\subsection{Correlation Is Not Causation}
\label{subsec:correlation-not-causation}
%%%%%%%%%%%%%%%%%%%%%%%%%%%%%%%%%%%%%%%%%%%%%%%%%%%%%%%%%%%%

A strong observed correlation does not prove that one variable causes
the other.

Possible explanations include:

\[
\boxed{
X\text{ causes }Y,
}
\]

\[
\boxed{
Y\text{ causes }X,
}
\]

\[
\boxed{
\text{a third variable influences both},
}
\]

or

\[
\boxed{
\text{the association arises through selection or chance}.
}
\]

Causal conclusions require stronger designs and assumptions.

%%%%%%%%%%%%%%%%%%%%%%%%%%%%%%%%%%%%%%%%%%%%%%%%%%%%%%%%%%%%
\subsection{Zero Correlation Does Not Mean No Relationship}
\label{subsec:zero-correlation-nonlinear}
%%%%%%%%%%%%%%%%%%%%%%%%%%%%%%%%%%%%%%%%%%%%%%%%%%%%%%%%%%%%

Correlation measures linear association.

A strong nonlinear relationship may have correlation near zero.

For example, if data follow approximately

\[
y=x^2
\]

with \(x\) symmetrically distributed around zero, positive and negative
linear tendencies may cancel.

Therefore a scatterplot remains essential.

%%%%%%%%%%%%%%%%%%%%%%%%%%%%%%%%%%%%%%%%%%%%%%%%%%%%%%%%%%%%
\subsection{Worked Example: Covariance and Correlation}
\label{subsec:worked-covariance-correlation}
%%%%%%%%%%%%%%%%%%%%%%%%%%%%%%%%%%%%%%%%%%%%%%%%%%%%%%%%%%%%

\begin{workedexamplebox}{Perfect Positive Association}

Suppose

\[
x=(1,2,3)
\]

and

\[
y=(2,4,6).
\]

Then

\[
y=2x.
\]

All points lie exactly on a positive-slope line.

Therefore,

\begin{answerbox}
\[
\boxed{
r=1.
}
\]
\end{answerbox}

\end{workedexamplebox}

%%%%%%%%%%%%%%%%%%%%%%%%%%%%%%%%%%%%%%%%%%%%%%%%%%%%%%%%%%%%
\subsection{Contingency Tables}
\label{subsec:contingency-tables}
%%%%%%%%%%%%%%%%%%%%%%%%%%%%%%%%%%%%%%%%%%%%%%%%%%%%%%%%%%%%

When two variables are categorical, their joint frequencies can be
organized into a contingency table.

For example,

\[
\begin{array}{c|cc|c}
&
B_1
&
B_2
&
\text{Total}
\\
\hline
A_1
&
n_{11}
&
n_{12}
&
n_{1\cdot}
\\
A_2
&
n_{21}
&
n_{22}
&
n_{2\cdot}
\\
\hline
\text{Total}
&
n_{\cdot1}
&
n_{\cdot2}
&
n
\end{array}
\]

The dot indicates summation across an index.

%%%%%%%%%%%%%%%%%%%%%%%%%%%%%%%%%%%%%%%%%%%%%%%%%%%%%%%%%%%%
\subsection{Marginal and Conditional Percentages}
\label{subsec:marginal-conditional-percentages}
%%%%%%%%%%%%%%%%%%%%%%%%%%%%%%%%%%%%%%%%%%%%%%%%%%%%%%%%%%%%

Marginal proportions describe one variable without conditioning on the
other.

For example,

\[
\boxed{
\frac{
n_{1\cdot}
}{n}
}
\]

is the overall proportion in category \(A_1\).

Conditional proportions divide by a row or column total.

For example,

\[
\boxed{
\frac{
n_{11}
}{
n_{1\cdot}
}
}
\]

is the proportion in \(B_1\) among observations already in \(A_1\).

%%%%%%%%%%%%%%%%%%%%%%%%%%%%%%%%%%%%%%%%%%%%%%%%%%%%%%%%%%%%
\subsection{Grouped Quantitative Data}
\label{subsec:grouped-quantitative-data}
%%%%%%%%%%%%%%%%%%%%%%%%%%%%%%%%%%%%%%%%%%%%%%%%%%%%%%%%%%%%

When only class intervals and frequencies are available, exact raw-data
statistics may be impossible to recover.

Suppose class \(i\) has midpoint

\[
m_i
\]

and frequency

\[
f_i.
\]

The grouped-data mean is approximated by

\[
\boxed{
\overline{x}
\approx
\frac{
\sum_i f_im_i
}{
\sum_i f_i
}.
}
\]

This treats every observation in a class as if it occurred at the class
midpoint.

%%%%%%%%%%%%%%%%%%%%%%%%%%%%%%%%%%%%%%%%%%%%%%%%%%%%%%%%%%%%
\subsection{Grouped Variance Approximation}
\label{subsec:grouped-variance}
%%%%%%%%%%%%%%%%%%%%%%%%%%%%%%%%%%%%%%%%%%%%%%%%%%%%%%%%%%%%

Using class midpoints,

\[
\boxed{
s^2
\approx
\frac{
\sum_i
f_i(m_i-\overline{x})^2
}{
n-1
}.
}
\]

The result is approximate because the exact positions of observations
inside each interval are unknown.

%%%%%%%%%%%%%%%%%%%%%%%%%%%%%%%%%%%%%%%%%%%%%%%%%%%%%%%%%%%%
\subsection{Geometric Mean}
\label{subsec:geometric-mean}
%%%%%%%%%%%%%%%%%%%%%%%%%%%%%%%%%%%%%%%%%%%%%%%%%%%%%%%%%%%%

For positive observations

\[
x_1,\ldots,x_n,
\]

the geometric mean is

\[
\boxed{
G
=
\left(
\prod_{i=1}^{n}x_i
\right)^{1/n}.
}
\]

Equivalently,

\[
\boxed{
\ln G
=
\frac1n
\sum_{i=1}^{n}\ln x_i.
}
\]

The geometric mean is useful for multiplicative growth.

%%%%%%%%%%%%%%%%%%%%%%%%%%%%%%%%%%%%%%%%%%%%%%%%%%%%%%%%%%%%
\subsection{Worked Example: Compound Growth}
\label{subsec:worked-geometric-mean}
%%%%%%%%%%%%%%%%%%%%%%%%%%%%%%%%%%%%%%%%%%%%%%%%%%%%%%%%%%%%

\begin{workedexamplebox}{Average Growth Factor}

Suppose yearly growth factors are

\[
1.10
\]

and

\[
1.20.
\]

The two-period combined factor is

\[
1.10(1.20)=1.32.
\]

The equivalent constant per-period growth factor is

\[
G
=
\sqrt{1.32}.
\]

\begin{answerbox}
\[
\boxed{
G
=
\sqrt{1.32}
\approx
1.149.
}
\]
\end{answerbox}

Thus the equivalent constant growth rate is approximately \(14.9\%\),
not the arithmetic mean \(15\%\) exactly.

\end{workedexamplebox}

%%%%%%%%%%%%%%%%%%%%%%%%%%%%%%%%%%%%%%%%%%%%%%%%%%%%%%%%%%%%
\subsection{Harmonic Mean}
\label{subsec:harmonic-mean}
%%%%%%%%%%%%%%%%%%%%%%%%%%%%%%%%%%%%%%%%%%%%%%%%%%%%%%%%%%%%

For nonzero observations,

\[
\boxed{
H
=
\frac{
n
}{
\sum_{i=1}^{n}
1/x_i
}.
}
\]

The harmonic mean is useful when averaging rates with a fixed numerator,
such as certain speed or ratio problems.

%%%%%%%%%%%%%%%%%%%%%%%%%%%%%%%%%%%%%%%%%%%%%%%%%%%%%%%%%%%%
\subsection{Arithmetic--Geometric--Harmonic Relationship}
\label{subsec:agh-means}
%%%%%%%%%%%%%%%%%%%%%%%%%%%%%%%%%%%%%%%%%%%%%%%%%%%%%%%%%%%%

For positive observations,

\[
\boxed{
H
\leq
G
\leq
\overline{x}.
}
\]

Equality holds when all observations are equal.

This provides a useful comparison among three different notions of
average.

%%%%%%%%%%%%%%%%%%%%%%%%%%%%%%%%%%%%%%%%%%%%%%%%%%%%%%%%%%%%
\subsection{Coefficient of Variation}
\label{subsec:coefficient-variation}
%%%%%%%%%%%%%%%%%%%%%%%%%%%%%%%%%%%%%%%%%%%%%%%%%%%%%%%%%%%%

For positive data with nonzero mean, the coefficient of variation is

\[
\boxed{
CV
=
\frac{s}{\overline{x}}.
}
\]

It is often expressed as a percentage:

\[
\boxed{
CV\%
=
100
\frac{s}{\overline{x}}.
}
\]

This measures variability relative to the mean.

%%%%%%%%%%%%%%%%%%%%%%%%%%%%%%%%%%%%%%%%%%%%%%%%%%%%%%%%%%%%
\subsection{When Coefficient of Variation Is Appropriate}
\label{subsec:cv-appropriate}
%%%%%%%%%%%%%%%%%%%%%%%%%%%%%%%%%%%%%%%%%%%%%%%%%%%%%%%%%%%%

The coefficient of variation is most meaningful for ratio-scale data
with a meaningful zero.

It is generally inappropriate when the mean is near zero or when the
measurement scale has an arbitrary zero.

%%%%%%%%%%%%%%%%%%%%%%%%%%%%%%%%%%%%%%%%%%%%%%%%%%%%%%%%%%%%
\subsection{Robust Statistics}
\label{subsec:robust-statistics-intro}
%%%%%%%%%%%%%%%%%%%%%%%%%%%%%%%%%%%%%%%%%%%%%%%%%%%%%%%%%%%%

A statistic is robust if its behavior is not excessively affected by a
small number of unusual observations or modest violations of model
assumptions.

Common robust summaries include:

\[
\boxed{
\text{median},
}
\]

\[
\boxed{
IQR,
}
\]

and

\[
\boxed{
\text{median absolute deviation}.
}
\]

%%%%%%%%%%%%%%%%%%%%%%%%%%%%%%%%%%%%%%%%%%%%%%%%%%%%%%%%%%%%
\subsection{Median Absolute Deviation}
\label{subsec:median-absolute-deviation}
%%%%%%%%%%%%%%%%%%%%%%%%%%%%%%%%%%%%%%%%%%%%%%%%%%%%%%%%%%%%

The median absolute deviation, abbreviated MAD, is

\[
\boxed{
MAD
=
\operatorname{median}
\left(
|x_i-\operatorname{median}(x)|
\right).
}
\]

This is a resistant measure of spread.

A rescaled MAD may be used as a robust estimator of normal-distribution
standard deviation.

%%%%%%%%%%%%%%%%%%%%%%%%%%%%%%%%%%%%%%%%%%%%%%%%%%%%%%%%%%%%
\subsection{Trimmed Means}
\label{subsec:trimmed-means}
%%%%%%%%%%%%%%%%%%%%%%%%%%%%%%%%%%%%%%%%%%%%%%%%%%%%%%%%%%%%

A trimmed mean removes a specified proportion of the smallest and
largest observations before computing the arithmetic mean.

For example, a \(10\%\) trimmed mean removes approximately the lowest
\(10\%\) and highest \(10\%\) of observations.

Trimmed means provide a compromise between:

\[
\boxed{
\text{efficiency of the mean}
}
\]

and

\[
\boxed{
\text{resistance of the median}.
}
\]

%%%%%%%%%%%%%%%%%%%%%%%%%%%%%%%%%%%%%%%%%%%%%%%%%%%%%%%%%%%%
\subsection{Winsorized Data}
\label{subsec:winsorized-data}
%%%%%%%%%%%%%%%%%%%%%%%%%%%%%%%%%%%%%%%%%%%%%%%%%%%%%%%%%%%%

Winsorization replaces extreme observations by selected boundary
quantiles rather than deleting them.

For example, values below a lower threshold may be replaced by that
threshold.

This limits the influence of extremes while retaining sample size.

%%%%%%%%%%%%%%%%%%%%%%%%%%%%%%%%%%%%%%%%%%%%%%%%%%%%%%%%%%%%
\subsection{Outliers Require Context}
\label{subsec:outliers-context}
%%%%%%%%%%%%%%%%%%%%%%%%%%%%%%%%%%%%%%%%%%%%%%%%%%%%%%%%%%%%

An outlier may represent:

\[
\boxed{
\text{a data-entry error},
}
\]

\[
\boxed{
\text{a measurement error},
}
\]

\[
\boxed{
\text{a rare but valid observation},
}
\]

or

\[
\boxed{
\text{evidence of a different population or mechanism}.
}
\]

Therefore an observation should not be removed merely because a
statistical rule labels it unusual.

%%%%%%%%%%%%%%%%%%%%%%%%%%%%%%%%%%%%%%%%%%%%%%%%%%%%%%%%%%%%
\subsection{Missing Data}
\label{subsec:missing-data-descriptive}
%%%%%%%%%%%%%%%%%%%%%%%%%%%%%%%%%%%%%%%%%%%%%%%%%%%%%%%%%%%%

Real datasets often contain missing observations.

Common missing-data codes include symbols such as

\[
NA
\]

or blank entries.

Missing values should not automatically be interpreted as zeros.

The mechanism causing missingness may affect statistical validity.

More formal missing-data theory belongs to later statistical modeling.

%%%%%%%%%%%%%%%%%%%%%%%%%%%%%%%%%%%%%%%%%%%%%%%%%%%%%%%%%%%%
\subsection{Data Cleaning and Descriptive Statistics}
\label{subsec:data-cleaning-descriptive}
%%%%%%%%%%%%%%%%%%%%%%%%%%%%%%%%%%%%%%%%%%%%%%%%%%%%%%%%%%%%

Before computing summaries, one should inspect:

\begin{itemize}
    \item impossible values;
    \item duplicate records;
    \item unit inconsistencies;
    \item coding errors;
    \item missing observations;
    \item unexpected categories;
    \item extreme values.
\end{itemize}

Descriptive statistics often serves as an early diagnostic tool for data
quality.

%%%%%%%%%%%%%%%%%%%%%%%%%%%%%%%%%%%%%%%%%%%%%%%%%%%%%%%%%%%%
\subsection{Comparing Groups}
\label{subsec:comparing-groups-descriptive}
%%%%%%%%%%%%%%%%%%%%%%%%%%%%%%%%%%%%%%%%%%%%%%%%%%%%%%%%%%%%

Suppose observations belong to groups

\[
A,
\quad
B,
\quad
C.
\]

A useful comparison should examine more than group means.

One should also compare:

\[
\boxed{
\text{sample size},
}
\]

\[
\boxed{
\text{spread},
}
\]

\[
\boxed{
\text{shape},
}
\]

and

\[
\boxed{
\text{outliers}.
}
\]

Boxplots and grouped histograms are common exploratory tools.

%%%%%%%%%%%%%%%%%%%%%%%%%%%%%%%%%%%%%%%%%%%%%%%%%%%%%%%%%%%%
\subsection{Standardized Group Comparisons}
\label{subsec:standardized-group-comparison}
%%%%%%%%%%%%%%%%%%%%%%%%%%%%%%%%%%%%%%%%%%%%%%%%%%%%%%%%%%%%

Suppose scores come from two different scales.

Raw values may not be directly comparable.

Standardization gives

\[
z
=
\frac{x-\mu}{\sigma}.
\]

This compares each observation relative to its own reference
distribution.

Thus \(z\)-scores are useful for comparing relative standing.

%%%%%%%%%%%%%%%%%%%%%%%%%%%%%%%%%%%%%%%%%%%%%%%%%%%%%%%%%%%%
\subsection{Descriptive Statistics and Probability Models}
\label{subsec:descriptive-probability-models}
%%%%%%%%%%%%%%%%%%%%%%%%%%%%%%%%%%%%%%%%%%%%%%%%%%%%%%%%%%%%

Sample statistics mirror population quantities from probability theory.

For example:

\[
\boxed{
\overline{x}
\longleftrightarrow
\mathbb{E}[X],
}
\]

\[
\boxed{
s^2
\longleftrightarrow
\operatorname{Var}(X),
}
\]

and

\[
\boxed{
r
\longleftrightarrow
\rho.
}
\]

This correspondence motivates statistical inference.

%%%%%%%%%%%%%%%%%%%%%%%%%%%%%%%%%%%%%%%%%%%%%%%%%%%%%%%%%%%%
\subsection{Empirical Versus Theoretical Distribution}
\label{subsec:empirical-versus-theoretical}
%%%%%%%%%%%%%%%%%%%%%%%%%%%%%%%%%%%%%%%%%%%%%%%%%%%%%%%%%%%%

A theoretical probability distribution gives probabilities such as

\[
F(x)
=
P(X\leq x).
\]

The empirical distribution gives

\[
F_n(x)
=
\frac1n
\sum_{i=1}^{n}
\mathbf{1}_{\{x_i\leq x\}}.
\]

Thus:

\[
\boxed{
F
=
\text{population distribution},
}
\]

while

\[
\boxed{
F_n
=
\text{sample distribution}.
}
\]

The Glivenko--Cantelli theorem later justifies uniform convergence of
\(F_n\) to \(F\) under i.i.d. sampling.

%%%%%%%%%%%%%%%%%%%%%%%%%%%%%%%%%%%%%%%%%%%%%%%%%%%%%%%%%%%%
\subsection{Descriptive Statistics Versus Statistical Inference}
\label{subsec:descriptive-versus-inference}
%%%%%%%%%%%%%%%%%%%%%%%%%%%%%%%%%%%%%%%%%%%%%%%%%%%%%%%%%%%%

Descriptive statistics asks:

\[
\boxed{
\text{What does this observed dataset look like?}
}
\]

Statistical inference asks:

\[
\boxed{
\text{What can this sample tell us about a larger population?}
}
\]

Inference requires probability models and sampling assumptions.

Descriptive summaries do not by themselves justify population-level
conclusions.

%%%%%%%%%%%%%%%%%%%%%%%%%%%%%%%%%%%%%%%%%%%%%%%%%%%%%%%%%%%%
\subsection{Exploratory Data Analysis}
\label{subsec:exploratory-data-analysis}
%%%%%%%%%%%%%%%%%%%%%%%%%%%%%%%%%%%%%%%%%%%%%%%%%%%%%%%%%%%%

Exploratory data analysis, abbreviated EDA, uses numerical summaries and
visualizations to discover patterns before formal modeling.

Typical EDA questions include:

\[
\text{Is the distribution skewed?}
\]

\[
\text{Are there several clusters?}
\]

\[
\text{Are there unusual observations?}
\]

\[
\text{Do two variables appear related?}
\]

\[
\text{Do groups differ visibly?}
\]

EDA is not a replacement for inference but an essential precursor to it.

%%%%%%%%%%%%%%%%%%%%%%%%%%%%%%%%%%%%%%%%%%%%%%%%%%%%%%%%%%%%
\subsection{Summary Statistics Table}
\label{subsec:summary-statistics-table}
%%%%%%%%%%%%%%%%%%%%%%%%%%%%%%%%%%%%%%%%%%%%%%%%%%%%%%%%%%%%

A common summary table for quantitative data includes:

\[
\begin{array}{c|c}
\text{Statistic} & \text{Meaning}\\
\hline
n & \text{sample size}\\
\overline{x} & \text{mean}\\
s & \text{standard deviation}\\
\min & \text{minimum}\\
Q_1 & \text{first quartile}\\
Q_2 & \text{median}\\
Q_3 & \text{third quartile}\\
\max & \text{maximum}\\
IQR & \text{interquartile range}
\end{array}
\]

No single row fully describes the dataset.

The statistics should be interpreted together.

%%%%%%%%%%%%%%%%%%%%%%%%%%%%%%%%%%%%%%%%%%%%%%%%%%%%%%%%%%%%
\subsection{Choosing Appropriate Summaries}
\label{subsec:choosing-summaries}
%%%%%%%%%%%%%%%%%%%%%%%%%%%%%%%%%%%%%%%%%%%%%%%%%%%%%%%%%%%%

For approximately symmetric quantitative data without severe outliers,
one often reports

\[
\boxed{
\text{mean and standard deviation}.
}
\]

For skewed data or data containing strong outliers, one often reports

\[
\boxed{
\text{median and IQR}.
}
\]

For categorical data, appropriate summaries include

\[
\boxed{
\text{counts and proportions}.
}
\]

%%%%%%%%%%%%%%%%%%%%%%%%%%%%%%%%%%%%%%%%%%%%%%%%%%%%%%%%%%%%
\subsection{Descriptive Statistics Workflow}
\label{subsec:descriptive-statistics-workflow}
%%%%%%%%%%%%%%%%%%%%%%%%%%%%%%%%%%%%%%%%%%%%%%%%%%%%%%%%%%%%

A practical descriptive-analysis workflow is:

\begin{enumerate}
    \item identify the observational units;
    \item identify each variable;
    \item classify variables as categorical or quantitative;
    \item inspect missing and impossible values;
    \item calculate sample size;
    \item produce frequency tables for categorical variables;
    \item graph quantitative distributions;
    \item calculate measures of center;
    \item calculate measures of spread;
    \item inspect quantiles and potential outliers;
    \item examine distribution shape;
    \item compare important groups;
    \item examine relationships among variables;
    \item choose summaries appropriate to the data structure;
    \item avoid making inferential claims without appropriate sampling
          assumptions.
\end{enumerate}

%%%%%%%%%%%%%%%%%%%%%%%%%%%%%%%%%%%%%%%%%%%%%%%%%%%%%%%%%%%%
\subsection{Common Errors}
\label{subsec:descriptive-statistics-common-errors}
%%%%%%%%%%%%%%%%%%%%%%%%%%%%%%%%%%%%%%%%%%%%%%%%%%%%%%%%%%%%

\subsubsection{Confusing Population and Sample}

A population is the full target collection.

A sample is the observed subset.

%%%%%%%%%%%%%%%%%%%%%%%%%%%%%%%%%%%%%%%%%%%%%%%%%%%%%%%%%%%%
\subsubsection{Confusing Parameter and Statistic}

A parameter describes a population.

A statistic is computed from sample data.

%%%%%%%%%%%%%%%%%%%%%%%%%%%%%%%%%%%%%%%%%%%%%%%%%%%%%%%%%%%%
\subsubsection{Averaging Category Codes}

If categories are coded as

\[
1,2,3,
\]

the numerical mean of those codes may have no meaningful interpretation
for nominal data.

%%%%%%%%%%%%%%%%%%%%%%%%%%%%%%%%%%%%%%%%%%%%%%%%%%%%%%%%%%%%
\subsubsection{Using Mean for Strongly Skewed Data Without Context}

The mean can be heavily influenced by a small number of extreme values.

Median and IQR may provide a more representative description.

%%%%%%%%%%%%%%%%%%%%%%%%%%%%%%%%%%%%%%%%%%%%%%%%%%%%%%%%%%%%
\subsubsection{Using \(n\) Instead of \(n-1\) for the Usual Sample Variance}

The common inferential sample variance is

\[
\boxed{
s^2
=
\frac1{n-1}
\sum
(x_i-\overline{x})^2.
}
\]

Using \(n\) produces the empirical second central moment instead.

Both quantities exist, but they serve different conventions and
purposes.

%%%%%%%%%%%%%%%%%%%%%%%%%%%%%%%%%%%%%%%%%%%%%%%%%%%%%%%%%%%%
\subsubsection{Reporting Variance in the Original Units}

Variance uses squared units.

Standard deviation uses the original units.

%%%%%%%%%%%%%%%%%%%%%%%%%%%%%%%%%%%%%%%%%%%%%%%%%%%%%%%%%%%%
\subsubsection{Calling Every Extreme Value an Error}

An outlier may be valid and scientifically important.

Statistical unusualness does not imply incorrectness.

%%%%%%%%%%%%%%%%%%%%%%%%%%%%%%%%%%%%%%%%%%%%%%%%%%%%%%%%%%%%
\subsubsection{Assuming Correlation Measures All Association}

Pearson correlation measures linear association.

Strong nonlinear relationships may have small correlation.

%%%%%%%%%%%%%%%%%%%%%%%%%%%%%%%%%%%%%%%%%%%%%%%%%%%%%%%%%%%%
\subsubsection{Assuming Correlation Implies Causation}

Association alone does not identify causal direction or rule out
confounding.

%%%%%%%%%%%%%%%%%%%%%%%%%%%%%%%%%%%%%%%%%%%%%%%%%%%%%%%%%%%%
\subsubsection{Comparing Raw Standard Deviations Across Different Scales}

A standard deviation depends on units.

Relative measures such as coefficient of variation may sometimes be
more useful when ratio-scale comparisons are meaningful.

%%%%%%%%%%%%%%%%%%%%%%%%%%%%%%%%%%%%%%%%%%%%%%%%%%%%%%%%%%%%
\subsubsection{Treating Histogram Shape as Unique}

Histogram appearance depends on bin width and bin boundaries.

Different reasonable choices can emphasize different features.

%%%%%%%%%%%%%%%%%%%%%%%%%%%%%%%%%%%%%%%%%%%%%%%%%%%%%%%%%%%%
\subsubsection{Ignoring Sample Size}

The same percentage may represent very different amounts of information
when based on \(n=10\) versus \(n=10{,}000\).

Always report or consider sample size.

%%%%%%%%%%%%%%%%%%%%%%%%%%%%%%%%%%%%%%%%%%%%%%%%%%%%%%%%%%%%
\subsubsection{Using Descriptive Statistics to Make Unsupported Population Claims}

A sample can be described accurately while still being unrepresentative
of the target population.

Sampling design matters.

%%%%%%%%%%%%%%%%%%%%%%%%%%%%%%%%%%%%%%%%%%%%%%%%%%%%%%%%%%%%
\subsection{Applications}
\label{subsec:descriptive-statistics-applications}
%%%%%%%%%%%%%%%%%%%%%%%%%%%%%%%%%%%%%%%%%%%%%%%%%%%%%%%%%%%%

\begin{applicationbox}

\textbf{Scientific Research.}
Descriptive statistics summarizes measurements, experimental outcomes,
and observational datasets before formal analysis.

\medskip

\textbf{Business.}
Means, medians, proportions, and variability summarize sales, customer
behavior, costs, and operational performance.

\medskip

\textbf{Engineering.}
Measurements of process output are summarized to detect variation and
possible quality problems.

\medskip

\textbf{Environmental Science.}
Temperature, rainfall, pollution, dissolved oxygen, and other
environmental measurements are summarized across time and location.

\medskip

\textbf{Education.}
Grades, attendance, survey responses, and program outcomes are summarized
using distributions, proportions, and measures of center and spread.

\medskip

\textbf{Healthcare Research.}
Patient measurements and outcomes are summarized before inferential
comparisons.

\medskip

\textbf{Data Science.}
Exploratory data analysis uses descriptive statistics to understand
datasets before modeling.

\medskip

\textbf{Government and Public Policy.}
Population characteristics, economic indicators, and survey data are
reported through statistical summaries and tables.

\end{applicationbox}

%%%%%%%%%%%%%%%%%%%%%%%%%%%%%%%%%%%%%%%%%%%%%%%%%%%%%%%%%%%%
\subsection{Exercises}
\label{subsec:descriptive-statistics-exercises}
%%%%%%%%%%%%%%%%%%%%%%%%%%%%%%%%%%%%%%%%%%%%%%%%%%%%%%%%%%%%

\begin{exercise}[1]
Define a population.
\end{exercise}

\begin{exercise}[1]
Define a sample.
\end{exercise}

\begin{exercise}[1]
Distinguish a parameter from a statistic.
\end{exercise}

\begin{exercise}[1]
Classify each variable as categorical or quantitative:
\begin{itemize}
    \item eye color;
    \item age;
    \item number of siblings;
    \item academic major;
    \item temperature.
\end{itemize}
\end{exercise}

\begin{exercise}[1]
Distinguish nominal and ordinal variables.
\end{exercise}

\begin{exercise}[1]
Distinguish discrete and continuous variables.
\end{exercise}

\begin{exercise}[1]
Define frequency and relative frequency.
\end{exercise}

\begin{exercise}[1]
Define the arithmetic mean.
\end{exercise}

\begin{exercise}[1]
Define the median.
\end{exercise}

\begin{exercise}[1]
Define the mode.
\end{exercise}

\begin{exercise}[1]
Define the sample variance.
\end{exercise}

\begin{exercise}[1]
Define the interquartile range.
\end{exercise}

\begin{exercise}[1]
Define a \(z\)-score.
\end{exercise}

\begin{exercise}[1]
Define the empirical distribution function.
\end{exercise}

\begin{exercise}[1]
Define sample covariance and sample correlation.
\end{exercise}

\begin{exercise}[2]
For the data

\[
3,\ 5,\ 5,\ 7,\ 10,
\]

compute the mean, median, mode, and range.
\end{exercise}

\begin{exercise}[2]
For the data

\[
2,\ 4,\ 6,
\]

compute the sample variance and standard deviation.
\end{exercise}

\begin{exercise}[2]
Show that the deviations from the sample mean sum to zero.
\end{exercise}

\begin{exercise}[2]
For the values

\[
10,\ 20,\ 30
\]

with weights

\[
0.2,\ 0.5,\ 0.3,
\]

compute the weighted mean.
\end{exercise}

\begin{exercise}[2]
Suppose

\[
Q_1=15
\]

and

\[
Q_3=27.
\]

Compute the IQR and the \(1.5IQR\) outlier fences.
\end{exercise}

\begin{exercise}[2]
Suppose

\[
\mu=100,
\qquad
\sigma=15.
\]

Find the \(z\)-score of

\[
x=130.
\]
\end{exercise}

\begin{exercise}[2]
What original value corresponds to

\[
z=-2
\]

when

\[
\mu=50,
\qquad
\sigma=5?
\]
\end{exercise}

\begin{exercise}[2]
For

\[
2,\ 4,\ 4,\ 6,\ 9,
\]

compute

\[
F_5(4).
\]
\end{exercise}

\begin{exercise}[2]
For a frequency table with frequencies

\[
3,\ 5,\ 2,
\]

compute the relative frequencies.
\end{exercise}

\begin{exercise}[2]
Find the geometric mean of

\[
4
\]

and

\[
9.
\]
\end{exercise}

\begin{exercise}[2]
Find the harmonic mean of

\[
3
\]

and

\[
6.
\]
\end{exercise}

\begin{exercise}[2]
Suppose

\[
\overline{x}=40,
\qquad
s=8.
\]

Compute the coefficient of variation.
\end{exercise}

\begin{exercise}[3]
Prove

\[
\sum_{i=1}^{n}
(x_i-\overline{x})
=
0.
\]
\end{exercise}

\begin{exercise}[3]
Derive the computational identity

\[
\sum
(x_i-\overline{x})^2
=
\sum x_i^2
-
n\overline{x}^2.
\]
\end{exercise}

\begin{exercise}[3]
Show that if

\[
y_i=ax_i+b,
\]

then

\[
\overline{y}
=
a\overline{x}+b.
\]
\end{exercise}

\begin{exercise}[3]
Show that under the same transformation,

\[
s_y^2
=
a^2s_x^2.
\]
\end{exercise}

\begin{exercise}[3]
Show that standardized sample values have mean zero.
\end{exercise}

\begin{exercise}[3]
Show that standardized sample values have sample standard deviation one.
\end{exercise}

\begin{exercise}[3]
For paired observations

\[
(1,2),
\quad
(2,4),
\quad
(3,6),
\]

compute covariance and correlation.
\end{exercise}

\begin{exercise}[3]
Construct a nonlinear dataset having a strong relationship but small or
zero Pearson correlation.
\end{exercise}

\begin{exercise}[3]
Explain why multiplying all \(x_i\) values by a positive constant does
not change correlation with \(y_i\).
\end{exercise}

\begin{exercise}[3]
Explain why multiplying all \(x_i\) values by a negative constant
reverses the sign of correlation.
\end{exercise}

\begin{exercise}[3]
Construct two datasets with the same mean but very different standard
deviations.
\end{exercise}

\begin{exercise}[3]
Construct two datasets with the same mean and standard deviation but
visibly different shapes.
\end{exercise}

\begin{exercise}[3]
Explain why mean and standard deviation alone cannot completely describe
a distribution.
\end{exercise}

\begin{exercise}[3]
Show that the empirical CDF is nondecreasing.
\end{exercise}

\begin{exercise}[3]
Show that the empirical CDF approaches \(0\) far below all observations
and \(1\) at or above the largest observation.
\end{exercise}

\begin{exercise}[3]
Explain the difference between a density histogram and a frequency
histogram.
\end{exercise}

\begin{exercise}[3]
Using class midpoints, derive the grouped-data mean approximation.
\end{exercise}

\begin{exercise}[3]
Explain why grouped-data variance is generally only approximate.
\end{exercise}

\begin{exercise}[3]
Compare the effect of a single extreme outlier on the mean, median,
standard deviation, and IQR.
\end{exercise}

\begin{exercise}[3]
Explain why the coefficient of variation is inappropriate for Celsius
temperature data.
\end{exercise}

\begin{exercise}[3]
Compare arithmetic and geometric means in a compound-growth setting.
\end{exercise}

\begin{exercise}[4]
Prove that the usual sample variance is an unbiased estimator of
population variance under i.i.d. sampling with finite variance.
\end{exercise}

\begin{exercise}[4]
Study different sample quantile conventions used by statistical
software.
\end{exercise}

\begin{exercise}[4]
Investigate robust estimators of location and scale.
\end{exercise}

\begin{exercise}[4]
Study the breakdown point of the mean and median.
\end{exercise}

\begin{exercise}[4]
Investigate influence functions for descriptive estimators.
\end{exercise}

\begin{exercise}[4]
Study robust covariance and correlation methods.
\end{exercise}

\begin{exercise}[4]
Investigate kernel density estimation as a smooth alternative to
histograms.
\end{exercise}

\begin{exercise}[4]
Study the Glivenko--Cantelli theorem for the empirical distribution
function.
\end{exercise}

\begin{exercise}[4]
Investigate the Dvoretzky--Kiefer--Wolfowitz inequality.
\end{exercise}

\begin{exercise}[4]
Study Q--Q plots and their use in comparing empirical and theoretical
distributions.
\end{exercise}

\begin{exercise}[4]
Investigate Anscombe's quartet and explain why numerical summaries alone
can be misleading.
\end{exercise}

\begin{exercise}[4]
Study modern visualization principles for high-dimensional datasets.
\end{exercise}

%%%%%%%%%%%%%%%%%%%%%%%%%%%%%%%%%%%%%%%%%%%%%%%%%%%%%%%%%%%%
\subsection{Conceptual Summary}
\label{subsec:descriptive-statistics-summary}
%%%%%%%%%%%%%%%%%%%%%%%%%%%%%%%%%%%%%%%%%%%%%%%%%%%%%%%%%%%%

A population is the full collection of interest.

A sample is an observed subset.

A population quantity is a parameter.

A sample-derived quantity is a statistic.

For quantitative observations

\[
x_1,\ldots,x_n,
\]

the sample mean is

\[
\boxed{
\overline{x}
=
\frac1n
\sum_{i=1}^{n}x_i.
}
\]

The sample variance is

\[
\boxed{
s^2
=
\frac1{n-1}
\sum_{i=1}^{n}
(x_i-\overline{x})^2.
}
\]

The sample standard deviation is

\[
\boxed{
s
=
\sqrt{s^2}.
}
\]

The interquartile range is

\[
\boxed{
IQR
=
Q_3-Q_1.
}
\]

A standardized score is

\[
\boxed{
z
=
\frac{x-\overline{x}}{s}
}
\]

or, using population values,

\[
\boxed{
z
=
\frac{x-\mu}{\sigma}.
}
\]

The empirical distribution function is

\[
\boxed{
F_n(x)
=
\frac1n
\sum_{i=1}^{n}
\mathbf{1}_{\{x_i\leq x\}}.
}
\]

For paired quantitative data, covariance is

\[
\boxed{
s_{xy}
=
\frac1{n-1}
\sum_{i=1}^{n}
(x_i-\overline{x})
(y_i-\overline{y}),
}
\]

and correlation is

\[
\boxed{
r
=
\frac{s_{xy}}{s_xs_y}.
}
\]

Descriptive statistics therefore summarizes data through

\[
\boxed{
\text{center}
+
\text{spread}
+
\text{shape}
+
\text{position}
+
\text{association}.
}
\]

These summaries provide the starting point for statistical inference.

%%%%%%%%%%%%%%%%%%%%%%%%%%%%%%%%%%%%%%%%%%%%%%%%%%%%%%%%%%%%
\subsection{Section Connections}
\label{subsec:descriptive-statistics-connections}
%%%%%%%%%%%%%%%%%%%%%%%%%%%%%%%%%%%%%%%%%%%%%%%%%%%%%%%%%%%%

\chapterconnection{
Descriptive Statistics begins Chapter 8 by converting the probability
concepts of Chapter 7 into empirical summaries calculated from observed
data. Sample means and variances mirror expectation and variance,
empirical distributions approximate probability distributions, and
sample covariance and correlation mirror their population counterparts.
The next section, Data Collection and Sampling, studies how observations
are obtained and whether a sample can legitimately represent a target
population. Later sections use descriptive statistics as the foundation
for sampling distributions, estimation, confidence intervals,
hypothesis testing, regression, experimental design, multivariate
statistics, Bayesian methods, time series, and statistical learning.
}

%%%%%%%%%%%%%%%%%%%%%%%%%%%%%%%%%%%%%%%%%%%%%%%%%%%%%%%%%%%%
% SECTION SOURCE TRACKING
%%%%%%%%%%%%%%%%%%%%%%%%%%%%%%%%%%%%%%%%%%%%%%%%%%%%%%%%%%%%

%------------------------------------------------------------
% 8.1 Descriptive Statistics
%------------------------------------------------------------
%
% Source basis:
%   - Standard introductory and mathematical statistics.
%   - Standard exploratory data-analysis concepts.
%
% Used for:
%   - populations and samples;
%   - parameters and statistics;
%   - variables and observations;
%   - categorical and quantitative data;
%   - nominal, ordinal, interval, and ratio scales;
%   - discrete and continuous variables;
%   - univariate, bivariate, and multivariate data;
%   - raw data;
%   - frequencies and relative frequencies;
%   - cumulative frequencies;
%   - bar charts;
%   - histograms;
%   - density histograms;
%   - measures of center;
%   - arithmetic and weighted means;
%   - median and mode;
%   - resistant statistics;
%   - range;
%   - population and sample variance;
%   - standard deviation;
%   - degrees of freedom;
%   - computational variance formulas;
%   - quartiles;
%   - percentiles;
%   - quantiles;
%   - IQR;
%   - five-number summaries;
%   - outlier fences;
%   - boxplots;
%   - z-scores;
%   - linear transformations;
%   - distribution shape;
%   - skewness;
%   - kurtosis;
%   - empirical distribution functions;
%   - scatterplots;
%   - covariance;
%   - correlation;
%   - contingency tables;
%   - grouped data;
%   - geometric and harmonic means;
%   - coefficient of variation;
%   - robust statistics;
%   - MAD;
%   - trimmed means;
%   - Winsorization;
%   - missing-data awareness;
%   - exploratory data analysis;
%   - applications and exercises.
%
% Notes:
%   - Data Collection and Sampling are developed next in Section 8.2.
%   - Probability counterparts of the main descriptive statistics were
%     developed in Chapter 7.
%   - Statistical inference begins formally with sampling distributions
%     and estimation in later sections.
%
%%%%%%%%%%%%%%%%%%%%%%%%%%%%%%%%%%%%%%%%%%%%%%%%%%%%%%%%%%%%